\documentclass[12pt]{beamer}


\mode<presentation>
{
	\usetheme{Warsaw}
	% or ...
	
	\setbeamercovered{transparent}
	% or whatever (possibly just delete it)
}


\usepackage[utf8]{inputenc}
\usepackage[T2A]{fontenc}
\usepackage[russian]{babel}
\usepackage{amsmath, amssymb}
\usepackage{booktabs}
\usepackage{graphicx}
\usepackage{multirow}

\setbeamertemplate{headline}{}
%\setbeamertemplate{footline}{}
\setbeamertemplate{footline}{
	\leavevmode%
	\hbox{%
		\begin{beamercolorbox}[wd=.5\paperwidth,ht=2.5ex,dp=1.125ex,leftskip=.3cm plus1fill,rightskip=.3cm]{author in head/foot}%
			\usebeamerfont{author in head/foot}\insertsection
		\end{beamercolorbox}%
		\begin{beamercolorbox}[wd=.5\paperwidth,ht=2.5ex,dp=1.125ex,leftskip=.3cm,rightskip=.3cm plus1fil]{title in head/foot}%
			\usebeamerfont{title in head/foot}\hfill\insertframenumber\,/\,\inserttotalframenumber
	\end{beamercolorbox}}%
	\vskip0pt%
}


\title[Об асимптотической мощности энергетического критерия]{Исследование асимптотической мощности энергетического критерия для проверки гипотезы о равенстве двух распределений в случае двух близких альтернатив}
\author{Сивков Вячеслав}
\institute{Санкт-Петербургский государственный университет}
\date{2025}

\begin{document}

\begin{frame}
    \titlepage
\end{frame}

\begin{frame}{Содержание}
    \tableofcontents
\end{frame}

\section{Введение}

\subsection{Постановка задачи}

\begin{frame}{Постановка задачи}
    \begin{block}{Основная гипотеза}
        Проверка гипотезы о равенстве двух распределений:
        \[
        H_0: F_1 = F_2 \quad \text{против} \quad H_1: F_1 \neq F_2
        \]
    \end{block}
    
    \begin{block}{Выборки}
        \begin{itemize}
            \item $X = (X_1, \ldots, X_n)$ из $F_1$
            \item $Y = (Y_1, \ldots, Y_m)$ из $F_2$
            \item Выборки независимые
        \end{itemize}
    \end{block}
\end{frame}

\begin{frame}{Альтернативная гипотеза}
    \begin{block}{Параметризация отличия}
        При $n = m$:
        \begin{align*}
        F_1(x) &= F(x) \\
        F_2(x) &= F\left(x(1 + \frac{h_2}{\sqrt{n}}) + \frac{h_1}{\sqrt{n}}\right)
        \end{align*}
    \end{block}
    
    \begin{itemize}
        \item $h_1$ - параметр отличия по сдвигу
        \item $h_2$ - параметр отличия по масштабу
        \item $F(x)$ - базовое распределение
    \end{itemize}
\end{frame}


\subsection{Цели исследования}
\begin{frame}{Цели исследования}
    \begin{enumerate}
        \item Изучение скорости сходимости эмпирических мощностей энергетического теста к асимптотической
        \item Сравнение мощностей пяти тестов:
        \begin{itemize}
            \item Энергетический тест (ET)
            \item Тест Андерсона-Дарлинга (AD)
            \item Тест Колмогорова-Смирнова (KS)
            \item Тест Крамера-фон Мизеса (CM)
            \item Тест Вилкоксона-Манна-Уитни (WMW)
        \end{itemize}
    	\item Сравнение мощностей различных версий энергетического теста
    \end{enumerate}
\end{frame}





\section{Методы проверки гипотез}

\subsection{ET}
\begin{frame}{Энергетический тест (ET)}
    \begin{block}{Статистика}
        \[
        \Phi_{nm} = \Phi_{AB} - \Phi_A - \Phi_B
        \]
        \begin{align*}
        \Phi_A &= \frac{1}{n^2} \sum_{1 \leq i \leq j \leq n} g(|X_i - X_j|) \\
        \Phi_B &= \frac{1}{m^2} \sum_{1 \leq i \leq j \leq m} g(|Y_i - Y_j|) \\
        \Phi_{AB} &= \frac{1}{nm} \sum_{i=1}^n \sum_{j=1}^m g(|X_i - Y_j|)
        \end{align*}
    \end{block}
    
    \begin{itemize}
        \item $g(x)$ - монотонная и непрерывная функция
        \item Используется статистика $n\Phi_{nn}$
    \end{itemize}
\end{frame}

\subsection{О вычислении асимптотической мощности энергетического теста}
\begin{frame}{О вычислении асимптотической мощности энергетического теста}
	тут будет теорема
\end{frame}


\subsection{AD}
\begin{frame}{Тест Андерсона-Дарлинга (AD)}
    \begin{block}{Статистика для k выборок}
        \[
        A_{kN}^2 = \frac{1}{N} \sum_{i=1}^k \frac{1}{n_i} \sum_{j=1}^{N-1} \frac{(NM_{ij} - jn_i)^2}{j(N-j)}
        \]
    \end{block}
    
    \begin{itemize}
        \item $N = n_1 + \ldots + n_k$ - общий объем выборок
        \item $M_{ij}$ - количество элементов i-й выборки $\leq Z_{(j)}$
        \item $Z_{(1)}, \ldots, Z_{(N)}$ - вариационный ряд
    \end{itemize}
\end{frame}

\subsection{KS}
\begin{frame}{Тест Колмогорова-Смирнова (KS)}
    \begin{block}{Статистика}
        \[
        KS = \max_{1 \leq i \leq N} |\hat{F}(Z_{(i)}) - \hat{G}(Z_{(i)})|
        \]
    \end{block}
    
    \begin{itemize}
        \item $\hat{F}$, $\hat{G}$ - эмпирические функции распределения
        \item $Z_{(1)}, \ldots, Z_{(N)}$ - объединенный вариационный ряд
        \item $N = n + m$
    \end{itemize}
\end{frame}


\subsection{CM}
\begin{frame}{Тест Крамера-фон Мизеса (CM)}
    \begin{block}{Статистика}
        \[
        CM = \frac{mn}{N^2} \sum_{i=1}^N (\hat{F}(Z_{(i)}) - \hat{G}(Z_{(i)}))^2
        \]
    \end{block}
    
    \begin{itemize}
        \item Основан на интеграле квадрата разности эмпирических функций
        \item Более чувствителен к различиям в хвостах распределений
    \end{itemize}
\end{frame}


\subsection{WMW}
\begin{frame}{Тест Вилкоксона-Манна-Уитни (WMW)}
    \begin{block}{Статистика}
        \[
        U_{nm} = \sum_{i=1}^n \sum_{j=1}^m \chi(X_i, Y_j)
        \]
        \[
        \chi(X_i, Y_j) = 
        \begin{cases}
        1, & X_i \geq Y_j \\
        0, & X_i < Y_j
        \end{cases}
        \]
    \end{block}
    
    \begin{itemize}
        \item Непараметрический тест
        \item Чувствителен к сдвигу распределений
    \end{itemize}
\end{frame}

\section{Моделирование}

\begin{frame}{Вычисление критического значения}
\end{frame}

\begin{frame}{Вычисление эмпирической мощности}
\end{frame}



\section{Результаты}

\subsection*{ET Нормальное отличие по масштабу}

\begin{frame}
табличка
\end{frame}


\subsection*{ET Коши отличие по сдвигу}
\begin{frame}
	табличка
\end{frame}

\subsection*{Все Нормальное сдвиг}
\begin{frame}
	график
\end{frame}

\subsection*{Все Коши масштаб}
\begin{frame}
	график
\end{frame}

\subsection*{Нормальное два ET}
\begin{frame}
	табличка
\end{frame}

\begin{frame}
	график
\end{frame}


\section{Выводы}

\begin{frame}{Выводы}
    \begin{enumerate}
        \item \textbf{Сходимость:} Эмпирические мощности быстро сходятся к асимптотическим для нормального распределения
        
        \item \textbf{Мощность тестов зависит от типа отличия:}
        \begin{itemize}
            \item Отличие по сдвигу: WMW наиболее мощный
            \item Отличие по масштабу: ET наиболее мощный
        \end{itemize}
        
        \item \textbf{Влияние распределения:}
        \begin{itemize}
            \item Для распределения Коши ET показывает лучшие результаты при отличиях по масштабу
            \item WMW не чувствителен к отличиям по масштабу
        \end{itemize}
    \end{enumerate}
\end{frame}

\begin{frame}{Заключение}
    \begin{block}{Выполненная работа}
        \begin{itemize}
            \item Проведено сравнение пяти статистических тестов
            \item Изучена сходимость эмпирических мощностей
            \item Проанализировано влияние типа отличия распределений
        \end{itemize}
    \end{block}
    
    \begin{block}{Практическая значимость}
        \begin{itemize}
            \item Рекомендации по выбору теста в зависимости от типа альтернативы
            \item Сравнение различных вариантов энергетического теста
        \end{itemize}
    \end{block}
\end{frame}

\section*{Литература}
\begin{frame}{Литература}
    \begin{thebibliography}{9}
    \bibitem{melas2024}
    В. Б. Мелас, Д. И. Сальников.
    \textit{Об асимптотической мощности «энергетического» теста для проверки гипотез о равенстве двух распределений}.
    Вестник СПбГУ. Математика. Механика. Астрономия, 2024.
    
    \bibitem{AD_article}
    F.-W. Scholz, M. A. Stephens.
    \textit{K-Sample Anderson–Darling Tests}.
    Journal of the American Statistical Association, 1987.
    
    \bibitem{WMW_article}
    H. B. Mann, D. R. Whitney.
    \textit{On a Test of Whether one of Two Random Variables is Stochastically Larger than the Other}.
    Annals of Mathematical Statistics, 1947.
    
    \bibitem{KS_CM_article}
    H. Buning.
    \textit{Kolmogorov-Smirnov and Cramér-von Mises Type Two-Sample Tests with Various Weight Functions}.
    Communications in Statistics - Simulation and Computation, 2001.
    \end{thebibliography}
\end{frame}

\section*{Fin}
\begin{frame}
    \centering
    \Huge{А}
    
    \vspace{1cm}
    \large{Вопросы?}
\end{frame}

\end{document}